\section{A hard problem: learning with errors}\label{lwe}

As we have seen (Section \ref{ec}), a lot of cryptography relies on hard math
problems. RSA is based on the difficulty of integer factorization;
elliptic curve cryptography depends on the discrete log assumption.

Our protocol for leveled FHE relies on a different hard problem: the
learning with errors problem (LWE). The problem is to solve systems of
linear equations, except that the equations are only approximately true
-- they permit a small ``error'' -- and instead of solving for rational
or real numbers, you are solving for integers modulo \(q\).

\subsection{A small example of an LWE problem}\label{lwe-small}

Here is a concrete example of an LWE problem and how one might attack it
``by hand.'' This exercise will make the inherent difficulty of the
problem quite intuitive.


\begin{problem}
We are working over \(\mathbb{F}_{11}\), and there is some secret vector
\(a = ( a_{1},\ldots,a_{4} )\). There are two sets of claims.
Each claim ``\(( x_{1},\ldots,x_{4} ):y\)'' purports the
relationship

\[y = a_{1}x_{1} + a_{2}x_{2} + a_{3}x_{3} + a_{4}x_{4} + \varepsilon,\mspace{6mu}\varepsilon \in \{ 0,1 \}.\]
(The \(\varepsilon\) is different from equation to equation.)

One of the sets of claims is ``genuine'' and comes from a consistent set
of \(a_{i}\), while the other set is ``fake'' and has randomly generated
\(y\) values. Tell them apart and find the correct secret vector
\(( a_{1},\ldots,a_{4} )\).

\begin{center}
\begin{tabular}{|c|c|}
\hline
Blue Set & Red Set \\
\hline\hline
(1, 0, 1, 7) : 2 & (5, 4, 5, 2) : 2 \\
\hline
(5, 8, 4, 10) : 2 & (7, 7, 7, 8) : 5 \\
\hline
(7, 7, 8, 5) : 3 & (6, 8, 2, 2) : 0 \\
\hline
(5, 1, 10, 6) : 10 & (10, 4, 4, 3) : 1 \\
\hline
(8, 0, 2, 4) : 9 & (1, 10, 8, 6) : 6 \\
\hline
(9, 3, 0, 6) : 9 & (2, 7, 7, 4) : 4 \\
\hline
(0, 6, 1, 6) : 9 & (8, 6, 6, 9) : 1 \\
\hline
(0, 4, 9, 7) : 5 & (10, 6, 1, 6) : 9 \\
\hline
(10, 7, 4, 10) : 10 & (3, 1, 10, 9) : 7 \\
\hline
(5, 5, 10, 6) : 9 & (2, 4, 10, 3) : 7 \\
\hline
(10, 7, 3, 1) : 9 & (10, 4, 6, 4) : 7 \\
\hline
(0, 2, 5, 5) : 6 & (8, 5, 7, 2) : 5 \\
\hline
(9, 10, 2, 1) : 3 & (4, 7, 0, 0) : 8 \\
\hline
(3, 7, 2, 1) : 6 & (0, 3, 0, 0) : 0 \\
\hline
(2, 3, 4, 5) : 3 & (8, 3, 2, 7) : 5 \\
\hline
(2, 1, 6, 9) : 3 & (4, 6, 6, 3) : 1 \\
\hline
\end{tabular}
\end{center}

\end{problem}

(\textbf{Solution sketch; can be skipped safely.}) One way to start
would be to define an \emph{information vector}
\[( x_{1},x_{2},x_{3},x_{4}|y|S ),\] where
\(S \subset {\mathbb{F}}_{11}\), to mean the statement
\[\sum a_{i}x_{i} = y + s,\text{ where }s \in S.\] In particular, a
purported approximation \(( x_{1},x_{2},x_{3},x_{4} ):y\) in
the LWE protocol corresponds to the information vector
\[( x_{1},x_{2},x_{3},x_{4}|y|\{ 0, - 1 \} ).\]
The benefit of this notation is that we can take linear combinations of
them. Specifically, if \(( X_{1}| y_{1} |S_{1} )\)
and \(( X_{2}| y_{2} |S_{2} )\) are information
vectors (where \(X_{i}\) are vectors), then

\[( \alpha X_{1} + \beta X_{2}| \alpha y_{1} + \beta y_{2} |\alpha S_{1} + \beta S_{2} ),\]

where \(\alpha S = \{ \alpha s|s \in S \}\) and
\(S + T = \{ s + t|s \in S,t \in T \}\).

We can observe the following:

\begin{enumerate}
\item
  If we obtain two vectors \(( X|y|S_{1} )\) and
  \(( X|y|S_{2} )\), then we have the information (assuming
  the vectors are accurate) \(( X|y|S_{1} \cap S_{2} )\). So
  if we are lucky enough, say, to have
  \(| S_{1} \cap S_{2} | = 1\), then we have found an exact
  equation with no error.
\item
  As we linearly combine vectors, their ``error part'' \(S\) gets bigger
  exponentially. So we can only add vectors very few times, ideally just
  one or two times, before they start being unusable.
\end{enumerate}

With these heuristics, we can start by looking at the Red Set, and make
vectors with many \(0\)'s in the same places.

\begin{enumerate}
\item
  Our eyes are drawn to the juicy-looking
  \(( 0,3,0,0|0|\{ 0, - 1 \} ),\) which
  immediately gives \(a_{2} \in \{ 0,7 \}\).
\item
  \(( 4,7,0,0|8|\{ 0, - 1 \} )\) gives
  \(4a_{1} + 7a_{2} \in \{ 7,8 \}.\) Also, since
  \(7a_{2} \in \{ 0,5 \},\)
  \[4a_{1} \in \{ 7,8 \} - \{ 0,5 \} = \{ 7,8,2,3 \},\]
  and \(a_{1} \in \{ 10,2,6,9 \}.\)
\item
  Adding
  \[( 10,4,4,3|1|\{ 0, - 1 \} ) + ( 7,7,7,8|5|\{ 0, - 1 \} )\]
  gives \(( 6,0,0,0|6|\{ 0, - 1, - 2 \} ),\) which
  is nice because it has 3 zeroes! This gives
  \(a_{1} \in \{ 1,8,10 \}.\) Combining with (2), we conclude
  that \(a_{1} = 10.\)
\item
  \ldots{}
\end{enumerate}

We omit the rest of the solution, which makes for some fun tinkering.

\subsection{General problem}

The LWE problem (Problem \ref{lwe}), like the discrete log
assumption, is one of those ``hard problems that you can build
cryptography on.'' The problem is to solve for constants
\[a_{1},\ldots,a_{n} \in {\mathbb{Z}}/q{\mathbb{Z}},\] given a bunch of
\textbf{approximate} equations of the form
\[y = a_{1}x_{1} + \cdots + a_{n}x_{n} + \epsilon,\] where each
\(\epsilon\) is a ``small'' error (for simplicity, say in \(\{ 0,1\}\)).
